% Part: first-order-logic
% Chapter: beyond
% Section: modal-logics

\documentclass[../../../include/open-logic-section]{subfiles}

\begin{document}

\olfileid{fol}{byd}{mod}

\olsection{Logica over wat kan en wat moet (modale logica)}

Bekijk deze voorwaardelijke zin:
\begin{quote}
  Als Jeremy alleen in die kamer is, dan is hij dronken en naakt en
  danst hij op de stoelen.
\end{quote}
Deze zin kan volgens de regels voor materiële implicatie waar zijn en toch
misleidend zijn. Hij lijkt namelijk te zeggen dat er méér verband is tussen de
voorwaarde en het gevolg dan alleen dit: de voorwaarde is onwaar, of het gevolg
is waar. De formulering wekt de indruk dat de zin niet alleen in deze ene wereld
waar is. In deze wereld kan hij al vanzelf waar zijn doordat Jeremy niet alleen
in de kamer is. Er lijkt ook te staan dat het gevolg \emph{waar zou zijn geweest}
\emph{als} de voorwaarde waar was geweest. Je kunt de zin dus zo lezen: hij is
niet alleen in deze wereld waar, maar in elke ``mogelijke'' wereld. Dan is hij
\emph{noodzakelijk} waar en niet toevallig alleen in deze wereld.

Modale logica is bedoeld om precies over zulke noodzakelijkheid te kunnen
redeneren. We krijgen modale propositielogica door aan de gewone
propositielogica een boxoperator toe te voegen. Als $!A$ !!a{formula} is, dan is
$\Box !A$ dus ook een formule. $\Box !A$ zegt intuïtief dat $!A$
\emph{noodzakelijk} waar is: dat is waar in elke mogelijke wereld. We gebruiken
$\Diamond !A$ meestal als afkorting van $\lnot \Box \lnot !A$. Die formule zegt
dat $!A$ \emph{mogelijk} waar is. We kunnen zulke modaliteit ook aan de
predicatenlogica toevoegen.

Met Kripke-!!{structure}s kunnen we een semantiek voor modale logica geven.
Kripke heeft deze semantiek oorspronkelijk zelfs met modale logica in gedachten
ontworpen. We beperken ons nu niet tot partiële ordeningen. Algemeen nemen we
een verzameling ``mogelijke werelden'' $P$ en een binaire
``toegankelijkheidsrelatie'' $\Atom{R}{x,y}$ tussen die werelden. De formule
$\Atom{R}{p,q}$ zegt intuïtief dat wereld~$q$ verenigbaar is met wereld~$p$.
Als we ons in wereld~$p$ ``bevinden'', moeten we dus nog rekening houden met de
mogelijkheid dat de wereld eruit had kunnen zien als wereld~$q$.

Modale logica heet soms ``intensionele'' logica, tegenover ``extensionele''
logica. De bedoelde semantiek van een extensionele logica, zoals de klassieke
logica, verwijst maar naar één wereld: de ``werkelijke'' wereld. De semantiek van
een intensionele logica gebruikt meer soorten objecten, waaronder verschillende
mogelijke werelden. Met modaliteit kunnen we niet alleen noodzakelijkheid
!!{structure}ren. We kunnen er ook andere talige constructies mee
!!{structure}ren. Daarbij geven we $\Box$ en $\Diamond$ per toepassing een
andere lezing. Bijvoorbeeld:
\begin{enumerate}
\item In bewijsbaarheidslogica lezen we $\Box !A$ als ``$!A$ is bewijsbaar''
  en $\Diamond !A$ als ``$!A$ is consistent.''
\item In epistemische logica, de logica over kennis en overtuigingen, kunnen we
  $\Box !A$ lezen als ``ik weet $!A$'' of ``ik geloof $!A$.''
\item In tijdslogica kunnen we $\Box !A$ lezen als ``$!A$ is altijd waar'' en
  $\Diamond !A$ als ``$!A$ is soms waar.''
\end{enumerate}

We willen de logica uitbreiden met regels en axioma's voor modaliteit. Het
systeem $\Log{S4}$ bestaat bijvoorbeeld uit de gewone axioma's en regels van de
propositielogica, plus deze axioma's:
\begin{align*}
& \Box (!A \lif !B) \lif (\Box !A \lif \Box !B)\\
& \Box !A \lif !A \\
& \Box !A \lif \Box \Box !A
\intertext{Daarnaast geldt de regel: ``concludeer uit $!A$ de formule
  $\Box !A$.'' Het systeem $\Log{S5}$ voegt daar dit axioma aan toe:}
& \Diamond !A \lif \Box \Diamond !A
\end{align*}
Voor verschillende toepassingen kunnen varianten van deze axioma's geschikt
zijn. Men neemt bijvoorbeeld meestal aan dat S5 het begrip logische
noodzakelijkheid weergeeft. Het handige is dat we gewoonlijk een semantiek
kunnen vinden waarvoor het !!{derivation} systeem correct en volledig is. Dat
doen we door natuurlijke eisen aan de toegankelijkheidsrelatie in de
Kripke-!!{structure}s te stellen. $\Log{S4}$ hoort bijvoorbeeld bij de klasse
van Kripke-!!{structure}s waarin de toegankelijkheidsrelatie reflexief en
transitief is. $\Log{S5}$ hoort bij de klasse van Kripke-!!{structure}s waarin
de toegankelijkheidsrelatie {\em universeel} is. Dat betekent dat elke wereld
vanuit elke andere wereld toegankelijk is. Daardoor geldt $\Box !A$ precies
dan wanneer $!A$ in elke wereld geldt.

\end{document}
